\chapter{Introducción}\label{cap:introduccion}

\section{Motivación}

Bogotá está en medio de la mayor transformación de su sistema de transporte desde la llegada de
TransMilenio en el año 2000. La Primera Línea del Metro tiene prevista su operación comercial para
marzo de 2028 \citep{emb2026}; la troncal de la avenida 68 está por entrar en servicio; el Regiotram
de Occidente abrirá por fases desde 2027 \citep{regiotram2026}; la Línea 2 está en licitación
\citep{emb2026l2}, y la Línea 3, con conexión a Soacha, está en estructuración \citep{emb2026l3}. Al
mismo tiempo, la ciudad sigue entre las más congestionadas del mundo: en 2025 la velocidad media en
hora pico fue de 18,9 km/h \citep{tomtom2026}, y la duración promedio de un viaje es de 48 minutos
\citep{sdm2024}.

Las decisiones que se tomen en los próximos años comprometen recursos del orden de decenas de
billones de pesos y su retorno depende de una variable que nadie conoce: cuánta gente querrá
desplazarse, desde dónde y hacia dónde, durante los próximos veinticinco años. La experiencia
internacional muestra que esa variable se pronostica mal. En la muestra de \citet{flyvbjerg2005}, el
pronóstico de pasajeros de los proyectos férreos sobreestimó la demanda real en más de 100\,\% en
promedio, y los costos de capital de esos proyectos excedieron lo presupuestado en 44,7\,\% en
promedio \citep{flyvbjerg2005wb}. Bogotá conoce de primera mano el otro lado de la incertidumbre: la
demanda troncal de TransMilenio cayó a la mitad en 2020 y en 2025 todavía no recuperaba el nivel de
2019 \citep{silla2026}.

Pese a ello, la evaluación tradicional de proyectos de transporte trabaja con un escenario central de
demanda y, a lo sumo, con un análisis de sensibilidad. Esa práctica tiene dos problemas. Primero, el
costo social es una función no lineal de la demanda (la congestión crece más que
proporcionalmente), de modo que el costo evaluado en la demanda esperada subestima el costo esperado
por la desigualdad de Jensen. Segundo, ignora el valor de la flexibilidad: un proyecto que puede
posponerse hasta conocer mejor la demanda vale más que el mismo proyecto comprometido hoy.

\section{Pregunta de investigación y objetivos}

La pregunta que guía el trabajo es:
\begin{quote}\itshape
¿Qué combinación de inversiones férreas y de política de suelo minimiza el costo social de la
movilidad en Bogotá y su región entre 2026 y 2050, cuando la demanda futura se modela como un
movimiento browniano geométrico y la decisión considera explícitamente el riesgo?
\end{quote}

La pregunta se responde en dos fases. La \emph{fase I} se sitúa en 2013, con la coyuntura económica y
de movilidad de ese año, cuando Bogotá contrató el diseño de un metro subterráneo, y evalúa la primera
línea con la información de entonces y con la que se conoció después. La \emph{fase II} actualiza el
modelo a las condiciones de 2025 y busca el diseño óptimo de la red hacia 2050. Con ambas fases se
responde una segunda pregunta, presente en el debate público:
\begin{quote}\itshape
¿Debió ser subterránea la Primera Línea del Metro, que hoy se construye elevada, dado el tamaño de
Bogotá? ¿Bajo qué condiciones lo sería?
\end{quote}

Los objetivos específicos son:
\begin{enumerate}[nosep]
  \item Formular un modelo estratégico de transporte de Bogotá-región, con demanda estocástica por zona
        y equilibrio modal con congestión, calibrado con datos públicos recientes.
  \item Simular la demanda con un MBG multivariado de solución exacta, calibrar su volatilidad con la
        serie de TransMilenio y estudiar la convergencia del estimador de Monte Carlo con técnicas de
        reducción de varianza y cuasi-Monte Carlo.
  \item Plantear el diseño urbano óptimo como un problema estocástico de selección de cartera de
        proyectos con criterios media--CVaR y restricción presupuestal, y resolverlo por enumeración
        sobre escenarios comunes.
  \item Reconstruir la coyuntura de 2013, recalibrar el modelo a ese año y evaluar \emph{ex ante} y
        \emph{ex post} las tipologías elevada y subterránea de la Primera Línea.
  \item Cuantificar el costo urbano del viaducto que haría preferible el subterráneo y compararlo con
        una estimación independiente, para la Primera Línea y para las líneas 2 y 3.
  \item Valorar la opción de esperar para construir la Línea 3 con mínimos cuadrados de Monte Carlo,
        validando el algoritmo contra la solución analítica de \citet{mcdonald1986}.
  \item Entregar el modelo como software abierto: un paquete de Python con pruebas, una aplicación en
        Streamlit y un repositorio público que permite reproducir cada cifra del documento.
\end{enumerate}

\section{Aportes}

El trabajo une dos tradiciones que rara vez se encuentran en la práctica de la planeación del
transporte en Colombia. De la matemática financiera toma los procesos estocásticos, la simulación de
Monte Carlo, las medidas coherentes de riesgo y la teoría de opciones reales; de la ingeniería de
transporte, la estructura de cuatro etapas simplificada, las funciones de congestión y el equilibrio
modal. El resultado no pretende sustituir el modelo de demanda de la Secretaría Distrital de
Movilidad, que trabaja con cientos de zonas y una red detallada; es un modelo \emph{estratégico}, en el
sentido de que su escala permite simular miles de futuros posibles y comparar decenas de diseños con
los mismos escenarios, algo inviable con un modelo detallado.

Hay tres resultados que el lector puede esperar. El primero es metodológico: la combinación de MBG con
números aleatorios comunes y cuasi-Monte Carlo permite ordenar diseños con errores de muestreo muy
pequeños frente a las diferencias entre ellos. El segundo es sustantivo: con los supuestos centrales,
la política de proximidad y el Regiotram del Norte dominan a las líneas de metro adicionales, y el
orden depende sobre todo de la tasa social de descuento y del valor del tiempo. El tercero es de
decisión: la Línea 3 tiene un valor de opción pequeño pero positivo, y la regla que lo captura es
esperar con una frontera explícita de ejercicio. El cuarto se refiere al metro: con la información y
la tasa de descuento de 2013 el elevado era preferible, pero la conclusión depende de la tasa de
descuento y del valor que se asigne al costo urbano del viaducto, y el documento identifica con
precisión los puntos de quiebre.

\section{Organización del documento}

La parte I reúne los fundamentos. El \cref{cap:contexto} resume el estado de la movilidad en Bogotá
y los proyectos en curso. El \cref{cap:marco} presenta el marco teórico: economía urbana estocástica,
el MBG, la simulación de Monte Carlo, las medidas de riesgo y las opciones reales. El
\cref{cap:modelo} formula el modelo de transporte y su calibración; el \cref{cap:simulacion}
desarrolla la simulación del MBG, y el \cref{cap:optimizacion} plantea el problema de diseño óptimo.
La parte II corresponde a la fase I: el \cref{cap:coyuntura2013} describe la coyuntura de 2013 y la
historia de la decisión del metro, y el \cref{cap:evaluacion2013} evalúa la primera línea con la
información de 2013 y con la observada. La parte III corresponde a la fase II: el
\cref{cap:resultados} presenta los resultados con las condiciones de 2025 y el \cref{cap:opciones}
valora la Línea 3 como opción real. La parte IV sintetiza: el \cref{cap:metro} reúne las conclusiones
sobre el metro elevado y el subterráneo, el \cref{cap:discusion} discute limitaciones e implicaciones y
el \cref{cap:conclusiones} concluye. Los apéndices recogen la notación y la descripción del software.
